%% kapa: exam=1 round=19 year=2008 period=1 form=A subject=accounting subject_ko=회계학 seq=34 number=114
%% source: Q-Net 감정평가사 기출문제 (cst00309) artlSeq=1001962 file=extracted/1008/1001962/2008년 제19회 감정평가사 1차시험 기출문제/1교시.hwp
%% conversion: kordoc 4.24.51 -> tools/gen1.py (auto); review=proofread-sonnet; flags=none
\begin{question}{114}
자본에 대한 다음 설명 중 가장 옳은 것은?
\begin{choices}
  \item 자본잉여금은 주식발행초과금, 감자차익 및 기타자본잉여금 등으로 자본금으로의 전입이나 결손금을 보전하는 경우를 제외하고는 다른 용도로 감소시킬 수 없다.
  \item 기타포괄손익누계액은 당기순이익의 결정에 반영되지 않은 항목이지만 자본거래 이외의 원천에 의하여 순자산의 변화를 초래하는 항목의 누적분이다.
  \item 기업의 영업활동, 투자활동 및 재무활동 중에서 이익잉여금은 영업활동과 재무활동의 결과 발생하는 이익의 누적분이다.
  \item 자기주식처분이익, 감자차익 및 주식매입선택권은 자본잉여금을 증가시킨다.
  \item 무상증자를 실시함에 따라 지분증권을 추가로 취득한 경우, 당해 지분증권을 자산으로 인식한다.
\end{choices}
\end{question}
\answer{2}
